\section{Conclusion: What the Bond Dimension Told Us}
\label{sec:conclusion}

% Closing section of the course: the single idea, the pipeline recap, a tally
% of where chi stayed small, grew or broke, and the open theoretical question.
% Every number in the table is stated, with its source, in Parts I-IV.

The course rested on a single idea. A physical field sampled on $N = 2^n$ grid
points is not an arbitrary vector: its well-separated scales are weakly coupled,
so it lies close to a manifold described by $\mathcal O(n\chi^2)$ numbers, and
the bond dimension $\chi$ measures how strongly the scales talk to each other.
Part~\partref{1} built that manifold and computed on it explicitly,
Part~\partref{2} learned to solve linear systems on it, Part~\partref{3}
stress-tested it on turbulence, and Part~\partref{4} showed that the same
structure tames the curse of dimensionality in data-driven models.

\paragraph{The pipeline, revisited.}
The \textbf{equation} was the heat equation in Parts~\partref{1}
and~\partref{2}, the incompressible Navier--Stokes equations in
Part~\partref{3}, and the unknown itself in Part~\partref{4}. The
\textbf{formulation} was finite differences on dyadic grids throughout, with a
projection method for the pressure in Part~\partref{3} and a discretization of
time as well as space in Part~\partref{4}. The \textbf{TN encoding} was TT-SVD
and the quantics map in Part~\partref{1}, TT-cross for functions, boundary
conditions and masks in Part~\partref{2}, and the space--time train of measured
or computed data in Part~\partref{4}. \textbf{Time evolution} was an MPO--MPS
product followed by rounding in Part~\partref{1}, a sweeping linear solve per
implicit step in Part~\partref{2} (\cref{subsec:linear-solvers}), the projection
loop of Part~\partref{3}, and the diagonal power $\Lambda^k$ of MPS-DMD in
Part~\partref{4}. \textbf{Solution retrieval} was point evaluation by
contracting the chain against the digits of an index in Part~\partref{1},
coarse-grained evaluation and pixel sampling in Part~\partref{3}, and modes,
eigenvalues and sparse equations in Part~\partref{4}. At no stage was the dense
$2^n$ vector formed.

\paragraph{Where $\chi$ stayed small, grew, or broke.}
The table collects the rank measurements stated in the notes. It is not a
benchmark; it is a map of where the single idea held.

\begin{center}
  \small
  \begin{tabular}{@{}p{0.13\textwidth}p{0.62\textwidth}p{0.12\textwidth}@{}}
    \toprule
    \textsc{$\chi$} & \textsc{Measurement} & \textsc{Where} \\
    \midrule
    stayed small
      & $e^{-3x}+\sin(4\pi x)$ at tolerance $10^{-10}$: $\chi = 3$ for every
        $n$ from 6 to 20. A straight edge: $\chi = 2$ from $N = 64$ to
        $N = 1024$. Space--time solves of the heat and wave equations:
        $\chi = 2$--$22$ on $1024\times1024$; Burgers: $\chi = 10$. Laminar
        cylinder and square flows at $\chi = 30$, the NACA~0040 airfoil at
        $\chi = 45$.
      & Parts~\partref{1}, \partref{3}, \partref{4} \\
    \addlinespace
    grew
      & A disc mask: $\chi = 17, 50, 128$ at $N = 64, 256, 1024$. Synthetic
        turbulence at tolerance $10^{-2}$: $\chi = 14, 26, 41, 68, 128$ for
        inertial ranges $k_{\max} = 4, 8, 16, 32, 64$, the proxy for growing
        $\mathrm{Re}$. A $1024^3$ DNS at $\mathrm{Re}_\lambda = 315$: $\chi =
        1000$ keeps the inertial range of $E(k)$, $\chi = 100$ runs the solver
        for 9--10 Kolmogorov time units, and below $\chi \approx 500$ even the
        interleaved encoding injects spurious small-scale fluctuations.
      & Parts~\partref{1}--\partref{3} \\
    \addlinespace
    broke
      & White noise saturates the space/time bond. A narrow Gaussian initial
        condition makes the space/time bond of the all-at-once heat solve grow.
        Without rounding, explicit heat stepping multiplies $\chi$ by $3$ per
        step. A noisy function under TT-cross: adding rank makes the error
        worse.
      & Parts~\partref{1}, \partref{2}, \partref{4} \\
    \bottomrule
  \end{tabular}
\end{center}

\noindent
Two readings follow. First, rank tracks scale coupling, not resolution or
sharpness: an aligned step costs $\chi = 2$ at any $N$, while a disc and a wide
inertial range cost rank because they couple every scale at once. Second,
\enquote{broke} always means \enquote{at the requested accuracy}: the noisy
standing wave of Part~\partref{4} goes from a saturated bond back to $\chi = 1$
when the tolerance is loosened from $10^{-10}$ to $10^{-2}$.

\paragraph{One structure, two problems.}
PDE solvers and data-driven models both describe dynamical systems, and the
curse of dimensionality hits both: the linear systems of implicit time stepping
on the one hand, the snapshot SVD of DMD and the dictionary of SINDy on the
other. The tensor train lifts it from both at once, through the same
$\mathcal O(n\chi^2)$ representation, and the space--time solve of
\cref{subsec:dmd-space-time} hands its output to MPS-DMD with no conversion in
between.

\paragraph{The comparison that matters.}
Against neural-network surrogates, the case for tensor networks does not rest on
raw accuracy. It rests on two things. Modes, frequencies, growth rates and
sparse equations come out of the method directly, not from a post-hoc probe.
And the rank is a single, monitorable complexity knob with a known meaning: the
number of retained singular values.

\paragraph{The open question.}
Every number in the table was measured, none was predicted. For ground states of
gapped one-dimensional Hamiltonians, the area law~\cite{Hastings2007} bounds
$\chi$ independently of system size, and for smooth functions rank bounds are
known~\cite{Lindsey2023}. For the solutions of the PDEs in this course there is
no comparable result. We have no sharp a-priori bound on $\chi$ for
Navier--Stokes as a function of $\mathrm{Re}$ (\cref{subsec:expressivity}), none
for the growth of $\chi$ over a long time integration, and none for the temporal
bond of measured data. Such a bound would turn the bet of this course into a
theorem, and it is the central open problem the notes leave behind.
